% GENERATED FILE. Edit canonical lessons, then run npm run generate:books.
% canonical-source-hash: fnv1a64:f164d88489bb16a5
% canonical-lessons: UR-C474-jarasim, UR-C474-mutaddi, UR-C474-infekshan, UR-C474-khasra, UR-C474-chechak

\chapter{Illnesses That Spread}
\label{ch:ur-illnesses-that-spread}
\hlchaptermodality{474}

\begin{chapteropening}
\textbf{By the end of this chapter:} \emph{I can talk about illnesses that spread.}

Complete the last lesson of chapter 474: I can talk about illnesses that spread.
\end{chapteropening}

\section[\texorpdfstring{\ur{جراثیم} (jarāsīm)}{جراثیم (jarāsīm)}]{\ur{جراثیم} (jarāsīm) — germs}

\label{lesson:UR-C474-jarasim}



\begin{quote}
Before the new one: say the Urdu for to limp, then the Urdu for an inflammation.
\end{quote}

\subsection*{You'll want to know: \ur{جراثیم}}
\textbf{\ur{جراثیم}} — \emph{jarāsīm} — \textquotedblleft{}germs\textquotedblright{}.

\begin{quote}
\ur{گندے} \ur{ہاتھوں} \ur{پر} \ur{جراثیم} \ur{ہوتے} \ur{ہیں۔}
\end{quote}

\begin{quote}
\emph{gande hāthoṅ par jarāsīm hote haiṅ.}
\end{quote}

\begin{quote}
There are germs on dirty hands.
\end{quote}

Masculine, used in the plural.

\begin{cousinweb}
\textbf{\ur{جراثیم}} (\emph{jarāsīm}) is usually traced, through Persian, to Arabic \emph{jarāthīm}, \textquotedblleft{}germs, seeds of a thing\textquotedblright{}.
\end{cousinweb}

\begin{grammarlens}[title={What you've built}]
One more word for describing what you lived and felt.
\end{grammarlens}

\subsection*{Your turn}
\begin{itemize}
\raggedright
  \item \emph{Say it:} \emph{jarāsīm}
  \item \emph{Say it:} \emph{jarāsīm}, once more
  \item \emph{Say it:} \emph{gande hāthoṅ par jarāsīm hote haiṅ.} — There are germs on dirty hands.
  \item \emph{Recall:} say the Urdu for to limp, then the Urdu for an inflammation, then say \emph{jarāsīm} again
\end{itemize}

\begin{tcolorbox}[breakable,colback=teal!4,colframe=teal!35!black,title={Before you move on}]
What does \ur{جراثیم} mean? (\textquotedblleft{}Germs\textquotedblright{}.) Say it once more.
\end{tcolorbox}
\section[\texorpdfstring{\ur{متعدی} (mutāddī)}{متعدی (mutāddī)}]{\ur{متعدی} (mutāddī) — contagious}

\label{lesson:UR-C474-mutaddi}



\begin{quote}
Before the new one: say the Urdu for an inflammation, then the Urdu for germs.
\end{quote}

\subsection*{You'll want to know: \ur{متعدی}}
\textbf{\ur{متعدی}} — \emph{mutāddī} — \textquotedblleft{}contagious\textquotedblright{}.

\begin{quote}
\ur{زکام} \ur{ایک} \ur{متعدی} \ur{بیماری} \ur{ہے۔}
\end{quote}

\begin{quote}
\emph{zukām ek mutāddī bīmārī hai.}
\end{quote}

\begin{quote}
A cold is a contagious illness.
\end{quote}

\begin{cousinweb}
\textbf{\ur{متعدی}} (\emph{mutāddī}) is usually traced, through Persian, to Arabic \emph{mutaʿaddī}, \textquotedblleft{}passing on to others\textquotedblright{}.
\end{cousinweb}

\begin{grammarlens}[title={What you've built}]
One more word for describing what you lived and felt.
\end{grammarlens}

\subsection*{Your turn}
\begin{itemize}
\raggedright
  \item \emph{Say it:} \emph{mutāddī}
  \item \emph{Say it:} \emph{mutāddī}, once more
  \item \emph{Say it:} \emph{zukām ek mutāddī bīmārī hai.} — A cold is a contagious illness.
  \item \emph{Recall:} say the Urdu for an inflammation, then the Urdu for germs, then say \emph{mutāddī} again
\end{itemize}

\begin{tcolorbox}[breakable,colback=teal!4,colframe=teal!35!black,title={Before you move on}]
What does \ur{متعدی} mean? (\textquotedblleft{}Contagious\textquotedblright{}.) Say it once more.
\end{tcolorbox}
\section[\texorpdfstring{\ur{انفیکشن} (infekshan)}{انفیکشن (infekshan)}]{\ur{انفیکشن} (infekshan) — an infection}

\label{lesson:UR-C474-infekshan}



\begin{quote}
Before the new one: say the Urdu for germs, then the Urdu for contagious.
\end{quote}

\subsection*{You'll want to know: \ur{انفیکشن}}
\textbf{\ur{انفیکشن}} — \emph{infekshan} — \textquotedblleft{}an infection\textquotedblright{}.

\begin{quote}
\ur{زخم} \ur{میں} \ur{انفیکشن} \ur{ہو} \ur{گیا۔}
\end{quote}

\begin{quote}
\emph{zakhm meṅ infekshan ho gayā.}
\end{quote}

\begin{quote}
The wound became infected.
\end{quote}

Masculine.

\begin{cousinweb}
\textbf{\ur{انفیکشن}} (\emph{infekshan}) is borrowed from English \emph{infection}.
\end{cousinweb}

\begin{grammarlens}[title={What you've built}]
One more word for describing what you lived and felt.
\end{grammarlens}

\subsection*{Your turn}
\begin{itemize}
\raggedright
  \item \emph{Say it:} \emph{infekshan}
  \item \emph{Say it:} \emph{infekshan}, once more
  \item \emph{Say it:} \emph{zakhm meṅ infekshan ho gayā.} — The wound became infected.
  \item \emph{Recall:} say the Urdu for germs, then the Urdu for contagious, then say \emph{infekshan} again
\end{itemize}

\begin{tcolorbox}[breakable,colback=teal!4,colframe=teal!35!black,title={Before you move on}]
What does \ur{انفیکشن} mean? (\textquotedblleft{}An infection\textquotedblright{}.) Say it once more.
\end{tcolorbox}
\section[\texorpdfstring{\ur{خسرہ} (khasrā)}{خسرہ (khasrā)}]{\ur{خسرہ} (khasrā) — measles}

\label{lesson:UR-C474-khasra}



\begin{quote}
Before the new one: say the Urdu for contagious, then the Urdu for an infection.
\end{quote}

\subsection*{You'll want to know: \ur{خسرہ}}
\textbf{\ur{خسرہ}} — \emph{khasrā} — \textquotedblleft{}measles\textquotedblright{}.

\begin{quote}
\ur{میرے} \ur{بھائی} \ur{کو} \ur{خسرہ} \ur{ہو} \ur{گیا۔}
\end{quote}

\begin{quote}
\emph{mere bhāī ko khasrā ho gayā.}
\end{quote}

\begin{quote}
My brother caught measles.
\end{quote}

Masculine.

\begin{cousinweb}
\textbf{\ur{خسرہ}} (\emph{khasrā}) is inherited: CDIAL lists the Hindi cousin \emph{khasrā}, \textquotedblleft{}scab, measles\textquotedblright{}, under Sanskrit \emph{khasa}, \textquotedblleft{}itch, scab\textquotedblright{} (CDIAL 3854): the illness is named for its rash.
\end{cousinweb}

\begin{grammarlens}[title={What you've built}]
One more word for describing what you lived and felt.
\end{grammarlens}

\subsection*{Your turn}
\begin{itemize}
\raggedright
  \item \emph{Say it:} \emph{khasrā}
  \item \emph{Say it:} \emph{khasrā}, once more
  \item \emph{Say it:} \emph{mere bhāī ko khasrā ho gayā.} — My brother caught measles.
  \item \emph{Recall:} say the Urdu for contagious, then the Urdu for an infection, then say \emph{khasrā} again
\end{itemize}

\begin{tcolorbox}[breakable,colback=teal!4,colframe=teal!35!black,title={Before you move on}]
What does \ur{خسرہ} mean? (\textquotedblleft{}Measles\textquotedblright{}.) Say it once more.
\end{tcolorbox}
\section[\texorpdfstring{\ur{چیچک} (chechak)}{چیچک (chechak)}]{\ur{چیچک} (chechak) — smallpox}

\label{lesson:UR-C474-chechak}



\begin{quote}
Before the new one: say the Urdu for an infection, then the Urdu for measles.
\end{quote}

\subsection*{You'll want to know: \ur{چیچک}}
\textbf{\ur{چیچک}} — \emph{chechak} — \textquotedblleft{}smallpox\textquotedblright{}.

\begin{quote}
\ur{ٹیکے} \ur{کی} \ur{وجہ} \ur{سے} \ur{اب} \ur{چیچک} \ur{نہیں} \ur{ہوتی۔}
\end{quote}

\begin{quote}
\emph{ṭīke kī vajah se ab chechak nahīṅ hotī.}
\end{quote}

\begin{quote}
Because of vaccination there is no smallpox now.
\end{quote}

Feminine.

\begin{cousinweb}
\textbf{\ur{چیچک}} (\emph{chechak}) is usually traced to Persian \emph{chēchak}, \textquotedblleft{}smallpox\textquotedblright{}.
\end{cousinweb}

\begin{grammarlens}[title={What you've built}]
One more word for describing what you lived and felt.
\end{grammarlens}

\subsection*{Your turn}
\begin{itemize}
\raggedright
  \item \emph{Say it:} \emph{chechak}
  \item \emph{Say it:} \emph{chechak}, once more
  \item \emph{Say it:} \emph{ṭīke kī vajah se ab chechak nahīṅ hotī.} — Because of vaccination there is no smallpox now.
  \item \emph{Recall:} say the Urdu for an infection, then the Urdu for measles, then say \emph{chechak} again
\end{itemize}

\begin{tcolorbox}[breakable,colback=teal!4,colframe=teal!35!black,title={Before you move on}]
What does \ur{چیچک} mean? (\textquotedblleft{}Smallpox\textquotedblright{}.) Say it once more.
\end{tcolorbox}
